\documentclass[UTF8,a4paper]{ctexart}
\usepackage{geometry,amsmath,amssymb,bm,multicol,array,xcolor}
\geometry{top=3.2mm,bottom=3.2mm,left=3.4mm,right=3.4mm}
\pagestyle{empty}
\setlength{\parindent}{0pt}
\setlength{\parskip}{0pt}
\setlength{\columnsep}{3.0mm}
\setlength{\columnseprule}{0.18pt}
\renewcommand{\arraystretch}{0.82}
\definecolor{B}{RGB}{0,78,158}
\definecolor{R}{RGB}{170,25,25}
\definecolor{G}{RGB}{0,105,65}
\newcommand{\pg}[1]{\begin{center}\colorbox{B}{\color{white}\bfseries #1}\end{center}\vspace{-2pt}}
\newcommand{\lab}[1]{\textbf{\textcolor{B}{#1}}}
\newcommand{\warn}[1]{\textbf{\textcolor{R}{#1}}}
\newcommand{\ok}[1]{\textbf{\textcolor{G}{#1}}}
\newcommand{\blk}[5]{\textbf{\textcolor{B}{#1}}\par
\lab{看}: #2\par
\lab{首}: #3\par
\lab{求}: #4\par
\lab{检}: #5\par\vspace{1pt}\hrule\vspace{1.2pt}}
\newcommand{\mini}[2]{\textbf{\textcolor{B}{#1}} #2\par\vspace{0.6pt}}
\newcommand{\chain}[3]{\lab{链}: #1 $\rightarrow$ #2 $\rightarrow$ #3\par}
\begin{document}
\fontsize{5.05}{5.36}\selectfont\raggedright

% ================= PAGE 1 =================
\pg{P1 总读题系统：先定模型，再套模板}
\begin{multicols}{3}
\blk{0. 全题通用四行答案}
{任何新题；特别是看不懂时。}
{第一行写模型句：“本题取\underline{\hspace{0.5cm}}为控制体/流线/截面/区域，属于\underline{\hspace{0.55cm}}模型。”}
{第二行写主方程；第三行写未知量链：已知量$\to$中间量$\to$目标量；第四行写数值与单位。}
{结尾必须查：单位；表压/绝压；方向正负；截面速度是否平均；能量是否跨泵/激波/损失。}

\blk{1. 30 秒圈题法}
{题干长、图复杂、符号多。}
{先圈五类词：物质(水/空气/理想/粘性)；对象(管/板/喷管/圆柱/闸门)；状态(定常/无旋/层流/Ma)；目标($p,Q,F,M,t$)；特殊词(泵/背压/激波/相似)。}
{按目标反推：求$p$找高度/能量/等熵；求$Q$找压差/面积/阻塞；求$F$找压力积分/CV/$C_D$；求$M$找压比/面积比/波表；求$t$找液面连续。}
{若同题出现粘性和Bernoulli，用“机械能+损失”；出现激波就分段，不能一式到底。}

\blk{2. 先判能不能用公式}
{你准备套 Bernoulli/H--P/等熵/势流/Moody 时。}
{先写适用条件。Bernoulli：定常、不可压、沿流线，实际流加损失/泵；H--P：圆管、层流、充分发展；等熵：绝热无摩擦无激波；势流：无粘无旋。}
{条件不满足就改公式：Bernoulli 加$h_L,H_p,H_T$；等熵遇激波要“波前等熵$\to$激波$\to$波后等熵”；层流/湍流先算$Re$。}
{不要把 $0.528$ 当万能压比；它只对应空气 $\gamma=1.4$、$M=1$临界状态。}

\blk{3. 截面编号法}
{图里有水箱、管路、泵、阀门、喷管、压力表。}
{第一步在图上标 1,2,3；每个截面写 $z,p,V,A$。大水面：$V\approx0$，开口表压$p=0$。}
{同一管段 $V=Q/A$；不同管径分开算。局部损失用元件所在管段速度。泵/水轮/激波/变径处必须分段。}
{表压只适合液体机械能；气体状态方程、汽蚀、蒸汽压必须用绝压。}

\blk{4. 单位换算急查}
{任何数值题。}
{先把所有量换 SI。}
{$1\,{\rm cm^2}=10^{-4}{\rm m^2}$；$Q({\rm m^3/s})=Q({\rm m^3/h})/3600$；$1{\rm kPa}=1000{\rm Pa}$；水头$m=p/(\rho g)$，水中$1{\rm kPa}\approx0.102m$。}
{所有最终答案写单位；力 N，压强 Pa/kPa，流量 $\rm m^3/s$，质量流量 kg/s。}

\blk{5. 结果异常回查}
{算出负压力、负损失、流量不守恒、Ma 支路怪、方向反。}
{先不重算，按五项查：单位、模型、边界、支路、符号。}
{负绝压：表压/绝压混；损失为负：上下游或速度头写反；质量流量不等：密度或面积错；Ma 两支：看亚声/超声支。}
{力题最后问“流体对管壁”还是“管壁对流体”；两者相反。}

\blk{6. 概念题三句模板}
{题问“简述/说明/为什么/定义/意义”。}
{第一句：按题干原词定义。}
{第二句：写物理原因或适用条件。第三句：写公式、判据或后果。}
{不写空泛话。每句都要含关键词：守恒、边界、相似、压缩、粘性、能量损失、总压变化等。}

\blk{7. 证明/推导题模板}
{题问证明、推导、由什么得到什么。}
{第一步写“在\underline{\hspace{0.45cm}}条件下，由\underline{\hspace{0.6cm}}方程出发”。}
{固定顺序：假设条件$\to$微分/积分方程$\to$边界条件/积分路径$\to$化简$\to$结论与限制。}
{最后必须写适用条件；否则推导题容易扣“理论前提”分。}

\blk{8. 查表题三行格式}
{题给 Moody、斜激波图、PM 表、阻力系数图、物性表。}
{先写“查表输入”。}
{输入：$Re,\varepsilon/d$查$\lambda$；$M_1$查正激波；$M_1,\theta$查斜激波；$\nu(M)$查PM；$Re$查$C_D$；温度查$\nu,\rho,p_v$。}
{查表后写“读得\underline{\hspace{0.5cm}}，代入\underline{\hspace{0.5cm}}”。不能只写数字。}

\end{multicols}\newpage

% ================= PAGE 2 =================
\pg{P2 静水压、测压、出流、应力：看到液面先用这页}
\begin{multicols}{3}
\blk{1. 静水压基本题}
{液体静止、闸门、水深、压强分布、压力中心。}
{首写 $p=p_0+\rho gh$，深度$h$从自由液面向下量。}
{平面合力 $F=\rho gh_C A$；压力中心 $y_D=y_C+I_C/(y_CA)$；曲面分解 $F_x=\rho gh_C A_x$，$F_y=\rho gV_p$。}
{压力方向垂直受压面；压力中心比形心深；开口液面表压为0，不是绝压为0。}

\blk{2. U 管/压差计}
{图有 U 形管、水银、油水多液体、两点压差。}
{从一个已知压强点出发，沿液柱“向下加$\rho gh$，向上减$\rho gh$”。}
{同一静止连通液体同高等压；测同一管流两点常用 $\Delta p=(\rho_m-\rho)g\Delta h$。}
{别把水银高度直接当水头；换算到目标流体压差。}

\blk{3. 小孔/水箱出流}
{大水箱、小孔、孔口、液面下降、求流量或时间。}
{首写孔口 Bernoulli：$V=C_v\sqrt{2gH}$，$Q=C_dA_o\sqrt{2gH}$。}
{定水头直接代；变水头接连续：$-A_t\,dH/dt=Q(H)$，再积分。}
{水箱大才取液面$V\approx0$；有收缩/流速系数按题给。}

\blk{4. 闸门/铰链/拉杆}
{图有矩形闸门、圆弧闸门、铰链、拉杆、求力。}
{先求水压力合力和作用点，不要先列铰链反力。}
{对铰链取矩求拉杆力；再用$x,y$力平衡求铰链反力。圆弧面用水平分力+竖直压力体。}
{矩臂要垂直距离；压力方向垂直面；表压大气部分一般相消。}

\blk{5. 应力张量题}
{题给应力张量$P$和平面法向$\bm n$，求应力矢量/法向/切向/夹角。}
{首写 $\bm p_n=P\bm n$。}
{法向分量 $p_N=\bm p_n\cdot\bm n$；切向矢量 $\bm p_T=\bm p_n-p_N\bm n$；夹角 $\cos\alpha=p_N/|\bm p_n|$。}
{$\bm n$必须是单位向量；法向分量不是矩阵对角元。}

\blk{6. 速度场判定}
{题给 $u,v,w$，问连续、有旋、势函数、流函数。}
{先写连续：不可压需 $\nabla\cdot\bm V=0$。}
{旋度 $\bm\Omega=\nabla\times\bm V$；无旋才有 $\phi$。二维不可压才有 $\psi$，$u=\psi_y,\ v=-\psi_x$。}
{先判条件再积分；不要直接求$\phi,\psi$。}

\blk{7. 给速度求 $\phi,\psi$}
{题问势函数/流函数。}
{先写先验条件：$\phi$需 $v_x-u_y=0$；$\psi$需 $u_x+v_y=0$。}
{$\phi_x=u$ 先对$x$积分，再用 $\phi_y=v$ 定未知函数；$\psi_y=u$ 先对$y$积分，再用 $-\psi_x=v$ 定函数。}
{常数可省；最后用偏导反查一遍。}

\blk{8. 非定常 U 管振荡}
{等截面 U 管、水柱振荡、不计粘性。}
{取液面位移$z$，写回复力 $2\rho gAz$。}
{质量 $\rho AL$，故 $\ddot z+(2g/L)z=0$，周期 $T=2\pi\sqrt{L/(2g)}$。}
{两液面位移方向相反；$L$ 是总液柱长度，不是一边长度。}

\blk{9. 物面边界条件}
{题问不穿透、无滑移、自由面、界面。}
{先翻译边界条件。}
{固壁无粘：$\bm V\cdot\bm n=0$；有粘再加 $\bm V=\bm V_w$；自由面 $p=p_a$；两流体界面速度连续、剪应力连续。}
{边界条件决定积分常数；漏边界条件推导很难满分。}

\end{multicols}\newpage

% ================= PAGE 3 =================
\pg{P3 Bernoulli、管路、泵、水轮、控制体受力}
\begin{multicols}{3}
\blk{1. 理想 Bernoulli}
{无粘、定常、不可压，沿流线；求$p,V,z$。}
{首写 $p/\rho g+V^2/2g+z=C$。}
{选两截面；大水面$V=0$；开口表压$p=0$；由能量差求速度或压强。}
{跨泵、水轮机、局部损失、长管摩擦时不能用纯 Bernoulli。}

\blk{2. 机械能方程}
{实际管路、泵、阀门、弯头、粗糙管、求$Q,p,H$。}
{首写 $p_1/\rho g+V_1^2/2g+z_1+H_p=p_2/\rho g+V_2^2/2g+z_2+h_L+H_T$。}
{$h_L=\sum\lambda L V^2/(2gd)+\sum\zeta V^2/(2g)$；$V=Q/A$；$\lambda$由$Re,\varepsilon/d$定。}
{泵加能放左边，水轮取能放右边；各损失用所在管段速度。}

\blk{3. 管路阻力流程}
{给$Q,d,L,\nu,\varepsilon,\zeta$，求压降或泵功。}
{首写 $V=Q/A,\ Re=Vd/\nu$。}
{层流 $\lambda=64/Re$；湍流查 Moody 或用近似；沿程 $h_f=\lambda(L/d)V^2/(2g)$；局部 $h_\zeta=\zeta V^2/(2g)$。}
{Moody 的$\lambda$通常是 Darcy 系数；不要和 Fanning 系数混。}

\blk{4. 已知压差求流量}
{水平管、水槽连管、给水头差求$Q$。}
{首写能量差 $\Delta H=h_L(Q)$。}
{先假设流态算$Q$；再算$Re$验证。湍流要迭代或查图；粗糙度由实际$Q$反推$\lambda$再反查$\varepsilon/d$。}
{最后检查假设流态是否自洽。}

\blk{5. 泵安装高度/汽蚀}
{吸水泵、要求最高安装高度、汽化压、安全余量。}
{从吸水池液面列到泵入口：$p_a/\rho g+z_0=p_B/\rho g+z_B+V^2/2g+h_L$。}
{限制 $p_{B,abs}\ge p_v+\Delta p_{\rm safe}$；或 $NPSH_a=(p_B-p_v)/\rho g+V^2/2g$。}
{汽蚀必须用绝压；水温给出时查$p_v,\nu$。}

\blk{6. 水轮机输出功率}
{管路中有水轮机，求输出功率。}
{首写机械能方程，把水轮机取走的水头记 $H_T$。}
{$H_T=H_{\rm up}-H_{\rm down}-h_L$；轴功率 $P=\eta\rho gQH_T$。}
{水轮机不是泵；符号反了会功率方向错。}

\blk{7. 控制体受力}
{弯管、喷嘴、射流冲板、法兰螺栓、支座反力。}
{先画 CV，标压力、重力、支座力、流入流出速度。}
{动量方程：$\sum\bm F=\sum_{\rm out}\rho Q\bm V-\sum_{\rm in}\rho Q\bm V$。分$x,y$写；压力力推入CV。}
{若题问“流体对管壁/螺栓”，最后取反。出口射大气用表压0。}

\blk{8. 喷嘴/弯管法兰}
{入口出口面积不同，求轴向力。}
{先用能量求 $p_1,p_2,V_1,V_2$。}
{再用 $x$向动量：$p_1A_1-p_2A_2+R_x=\rho Q(V_2-V_1)$。}
{压力用表压；$R_x$是外界对流体，题问法兰受力取相反。}

\blk{9. 动量矩/旋转机械}
{水轮机、泵叶轮、喷流打转轮，求转矩或功率。}
{首写角动量方程。}
{$\sum M_O=\sum_{\rm out}\rho Q(rV_\theta)-\sum_{\rm in}\rho Q(rV_\theta)$；功率 $P=M\omega$。}
{只取切向速度分量；半径和旋转方向决定正负。}

\end{multicols}\newpage

% ================= PAGE 4 =================
\pg{P4 粘性流动、壁律、边界层、外绕阻力}
\begin{multicols}{3}
\blk{1. 圆管层流 H--P}
{圆管、层流、充分发展、求速度分布/流量/压降。}
{首写 $Re<2300$ 且充分发展。}
{$u(r)=\frac{-dp/dx}{4\mu}(R^2-r^2)$；$\bar u=u_{\max}/2$；$Q=\pi R^4\Delta p/(8\mu L)$；$\lambda=64/Re$。}
{H--P 不能用于湍流和入口发展段。}

\blk{2. 平板 Couette}
{两平板间流体，上板速度$U$，无压差。}
{首写 $\mu d^2u/dy^2=0$。}
{边界 $u(0)=0,u(h)=U$；得 $u=Uy/h$，剪应力 $\tau=\mu U/h$。}
{方向：壁面对流体剪力与相对运动相反；题问板受力取反。}

\blk{3. Couette + Poiseuille}
{两板间有上板运动，又有压强梯度/重力沿流向。}
{首写 $\mu u''=dp/dx-\rho g_x$。}
{积分两次，用两个壁面无滑移定常数；流量 $Q=b\int_0^h udy$。}
{若中点速度为零，用它反求$U$或压强梯度。}

\blk{4. 两层流体 Couette}
{两种粘度流体分层，上板运动，下板固定。}
{首写每层 $u_i=A_i y+B_i$。}
{边界：两壁无滑移；界面速度连续；界面剪应力连续 $\mu_1u_1'=\mu_2u_2'$。}
{不要令两层速度梯度相等；相等的是剪应力。}

\blk{5. 湍流壁律/边界层厚度}
{题给壁面剪切、$u_\tau,y^+$、壁律、求边界层厚度。}
{首写 $u_\tau=\sqrt{\tau_w/\rho}$，$y^+=yu_\tau/\nu$，$u^+=u/u_\tau$。}
{粘性底层 $u^+=y^+$；对数层 $u^+=2.5\ln y^++5.5$；边界层外缘常取$y^+\approx30$或$u\approx U$。}
{$\bar u$ 是截面平均速度，不等于局部$u(y)$。}

\blk{6. 平板边界层阻力}
{零攻角平板、给$U,L,b,\nu$，求阻力/厚度。}
{首写 $Re_L=UL/\nu$。}
{层流：$\delta=5L/\sqrt{Re_L}$，$\bar C_f=1.328/\sqrt{Re_L}$；湍流常用 $\bar C_f=0.074/Re_L^{1/5}$。$D=\bar C_f(0.5\rho U^2)A_{wet}$。}
{单面$A=Lb$，双面$2Lb$；若给$Re_{cr}$，先算$x_{cr}=Re_{cr}\nu/U$分段。}

\blk{7. 边界层积分题}
{给速度剖面 $u/U=f(y/\delta)$，求$\delta,\theta,C_D$。}
{设 $\eta=y/\delta$。}
{$\delta^*=\delta\int_0^1(1-f)d\eta$；$\theta=\delta\int_0^1 f(1-f)d\eta$；$\tau_w=\mu U f'(0)/\delta$；零压梯度 $\tau_w=\rho U^2d\theta/dx$。}
{最后阻力：单面 $C_D=2\theta(L)/L$。}

\blk{8. 分离判断}
{问边界层分离、逆压梯度、尾流、阻力危机。}
{首写逆压梯度 $dp/dx>0$ 使近壁流体减速。}
{分离点条件 $\tau_w=\mu(\partial u/\partial y)_0=0$；之后近壁可能回流。湍流抗逆压梯度强，分离后移，阻力危机时$C_D$骤降。}
{理想势流阻力为0不能解释真实阻力；真实阻力看粘性和分离。}

\blk{9. 外绕阻力}
{球、圆柱、汽车、阻力系数图。}
{首写 $Re=UL/\nu$，查 $C_D(Re)$。}
{$D=C_D(0.5\rho U^2)A_{\rm front}$；终速题令重力差=阻力。低$Re$小球可用 Stokes：$D=3\pi\mu dU$。}
{外绕阻力面积通常是迎风面积，不是湿面积；平板摩擦才用湿面积。}

\end{multicols}\newpage

% ================= PAGE 5 =================
\pg{P5 势流叠加、圆柱、镜像、升力}
\begin{multicols}{3}
\blk{1. 势流总判断}
{题写理想、不可压、无旋、势函数、流函数、叠加。}
{首写 $\nabla\times\bm V=0\Rightarrow\bm V=\nabla\phi$；不可压 $\nabla^2\phi=0$。二维不可压：$u=\psi_y,v=-\psi_x$。}
{复杂流场直接叠加：$\phi=\sum\phi_i,\ \psi=\sum\psi_i$；速度由偏导求。}
{Bernoulli 可在无旋势流全场任意两点用；真实粘性阻力不能用势流给。}

\blk{2. 基元流识别表}
{看见源、汇、点涡、偶极子、均匀流。}
{首写对应$\phi,\psi$。}
{均匀流：$\phi=Ur\cos\theta,\psi=Ur\sin\theta$；源：$v_r=Q/(2\pi r),\psi=Q\theta/(2\pi)$；涡：$v_\theta=\Gamma/(2\pi r)$；偶极：$\phi=M\cos\theta/r$。}
{半平面/壁面问题常把$2\pi$换成$\pi$，先看是否有镜像。}

\blk{3. 源+均匀流半体}
{均匀流中有点源，求驻点/分界线/半体形状。}
{首写总流函数 $\psi=Ur\sin\theta+Q\theta/(2\pi)$。}
{驻点：$v_r=0,v_\theta=0$，得 $\theta=\pi,\ r=Q/(2\pi U)$；分界流线过驻点，$\psi=Q/2$。}
{远处半体宽度 $b=Q/U$；方向由源/汇号判。}

\blk{4. 圆柱绕流无环量}
{均匀流绕圆柱、圆柱压力、驻点、达朗贝尔佯谬。}
{首写“均匀流+偶极子叠加，圆柱面$r=a$为$\psi=0$流线”。}
{$\phi=U(r+a^2/r)\cos\theta$，$\psi=U(r-a^2/r)\sin\theta$；$v_r=U(1-a^2/r^2)\cos\theta$；$v_\theta=-U(1+a^2/r^2)\sin\theta$。}
{壁面$C_p=1-4\sin^2\theta$；压力积分阻力升力均为0，只适合理想无粘。}

\blk{5. 圆柱有环量/升力}
{圆柱叠加点涡、机翼升力、Kutta-Joukowski。}
{首写壁面切向速度 $v_\theta=-2U\sin\theta+\Gamma/(2\pi a)$。}
{用 Bernoulli 得压力分布；升力 $L'=\rho U\Gamma$。驻点由 $v_\theta=0$，即 $\sin\theta=\Gamma/(4\pi Ua)$。}
{升力方向按来流和环量右手关系；阻力仍为0是理想结论。}

\blk{6. 镜像法}
{点源/点涡靠近直壁、半平面、壁面吸水。}
{首写“固壁是流线，法向速度为0”。}
{源/汇靠直壁用同号镜像；点涡靠直壁用反号镜像；写总$\psi$后令壁面$\psi=C$，驻点由速度为0。}
{半平面源汇常见分母从$2\pi$变$\pi$；符号错会导致壁面有穿透速度。}

\blk{7. 边角绕流}
{楔角、角域、问速度在角点是否无穷。}
{首写 $\phi=Ar^n\cos n\theta,\psi=Ar^n\sin n\theta$。}
{壁面为$\psi=0$，夹角$\alpha$满足 $n=\pi/\alpha$；速度 $V\sim r^{n-1}$。}
{$\alpha<\pi$速度趋0；$\alpha>\pi$速度趋无穷，实际要靠粘性修正。}

\blk{8. 非定常 Bernoulli}
{无旋非定常、U管振荡、势函数含$t$。}
{首写 $\partial\phi/\partial t+V^2/2+p/\rho+\Pi=f(t)$。}
{同一时刻两点相减消去$f(t)$；若定常则$\partial\phi/\partial t=0$。}
{只适用于无旋；有损失/粘性要改机械能方程。}

\blk{9. Kelvin/Lagrange/Helmholtz}
{证明环量守恒、无旋保持、涡线涡管。}
{首写条件：理想、正压、保守体力。}
{Kelvin：$D\Gamma/Dt=0$；Lagrange：若初始$\Omega=0$，以后仍无旋；Helmholtz：涡线随体、涡管强度守恒、涡线不能在流体内开始/终止。}
{必须写条件；不满足粘性/非正压时不能直接用。}

\end{multicols}\newpage

% ================= PAGE 6 =================
\pg{P6 可压缩、喷管、激波、膨胀波、相似}
\begin{multicols}{3}
\blk{1. 可压缩入门}
{空气、声速、马赫数、压缩性、$Ma$。}
{首写 $Ma=V/a,\ a=\sqrt{\gamma RT}$。}
{$Ma\le0.3$ 常可忽略压缩；否则用可压关系。等熵：$T_0/T=1+(\gamma-1)M^2/2$，$p_0/p=(T_0/T)^{\gamma/(\gamma-1)}$。}
{气体状态方程 $p=\rho RT$ 必须用绝压和绝温 K。}

\blk{2. Laval 喷管判型}
{收缩/缩扩喷管、喉部、背压、阻塞、最大流量。}
{首写比较背压与临界压比。空气临界 $p^*/p_0=0.528,\ T^*/T_0=0.833$。}
{若阻塞：喉部$M=1,A_t=A^*$，$\dot m_{\max}=0.0404A^*p_0/\sqrt{T_0}$。未阻塞：出口内压常等于背压，用等熵压比求$M$。}
{$p_b$是外部背压，不一定等于出口内压$p_e$；阻塞后降背压不增大质量流量。}

\blk{3. 面积--马赫数}
{给 $A/A^*$ 求 $M$ 或给 $M$ 求面积。}
{首写面积关系，并标明有亚声/超声两支。}
{$A/A^*=1$ 时 $M=1$；收缩段亚声加速，扩张段若已超声继续加速。}
{支路由题意决定：喷管入口通常亚声，喉后缩扩喷管可能超声。}

\blk{4. 正激波}
{题写正激波、管内激波、波前波后 1$\to$2。}
{首写“波前等熵求$M_1,p_1,T_1$，过正激波求$M_2,p_2,T_2,p_{02}$”。}
{空气常用：$M_2^2=(1+0.2M_1^2)/(1.4M_1^2-0.2)$；$p_2/p_1=1+\frac{2\gamma}{\gamma+1}(M_1^2-1)$；再用波后等熵。}
{正激波后 $T_0$不变、$p_0$下降；不能继续用波前总压。}

\blk{5. 斜激波}
{超声速遇压缩角、楔形、凹角、斜激波角$\beta$。}
{首写 $M_{n1}=M_1\sin\beta$，法向按正激波。}
{查/解 $\theta$-$\beta$-$M$；弱解常见，强解损失大。波后切向速度不变，法向按正激波变。}
{$\beta=\mu=\arcsin(1/M_1)$ 为马赫波；$\beta=90^\circ$ 为正激波；超过最大折转角会脱体。}

\blk{6. Prandtl--Meyer 膨胀波}
{超声速遇凸角、喷管外膨胀、低压区。}
{首写 $\nu(M_2)=\nu(M_1)+\delta$。}
{查 PM 函数得$M_2$；再用等熵关系求$p_2,T_2,\rho_2$。膨胀波使速度升高、压强温度降低。}
{PM 是连续等熵膨胀，$p_0$不变；不要用激波关系。}

\blk{7. 管外膨胀/压缩波系}
{喷管出口压力和背压不匹配，图有外部波。}
{先由$A_e/A_t$等熵求喷管出口内压$p_e$。}
{若$p_e>p_b$：出口后继续膨胀，PM；若$p_e<p_b$：外部压缩，斜激波/正激波；若$p_e=p_b$：完全膨胀。}
{只在匹配或全亚声特殊情况才可令$p_e=p_b$。}

\blk{8. Fanno 等截面摩擦管}
{可压气体在等截面长管，有摩擦、问最大长度/临界。}
{首写：等截面、绝热、有摩擦，查 Fanno。}
{由入口$M_i$查 $4fL^*/D$；实际 $L$ 与 $L^*$ 比较。摩擦使流动趋向 $M=1$，$T_0$不变，$p_0$下降。}
{这不是等熵；不能用喷管面积关系替代。}

\blk{9. 波阻/声障/声爆概念}
{题问波阻、声障、声爆、激波熵增。}
{首句：超声速流经激波发生不可逆压缩，熵增导致总压损失。}
{波阻来自激波造成的动量损失和压力分布变化；正激波损失最大。声爆是激波传播到地面形成突跃压力。}
{概念题要写“与摩擦不同，属于压缩波不可逆损失”。}

\blk{10. 量纲分析}
{题给多个物理量，要求建立无量纲关系。}
{首写 Buckingham：$n$个量、$r$个基本量纲，得$n-r$个$\pi$。}
{选重复变量：量纲独立、不含待求量、覆盖$M,L,T$；设 $\pi=X\rho^aV^bL^c$ 配平指数。}
{每个$\pi$必须无量纲；重复变量不能全来自同一量纲。}

\blk{11. 模型相似}
{模型实验、缩尺、相似准则、求原型量。}
{先写几何相似，再选主导准则。}
{$Re=\rho VL/\mu$ 主粘性；$Fr=V/\sqrt{gL}$ 主重力；$Ma$ 主压缩；$We$ 主表面张力；$Eu=\Delta p/(\rho V^2)$ 主压强。}
{同流体缩尺常 $Re$ 与 $Fr$ 冲突；保主导准则叫部分相似。}

\blk{12. 水击/水波/明渠}
{题问水击波速、深水波、扰动能否上游传播。}
{水击首写 $\Delta p=\rho a\Delta V$；深水波首写 $h>\lambda/2$ 或 $kh>\pi$；明渠首写 $Fr=V/\sqrt{gh}$。}
{水击 $a\approx\sqrt{K/\rho}/\sqrt{1+KD/(Ee)}$；深水 $c=\sqrt{g\lambda/(2\pi)}$；浅水 $c=\sqrt{gh}$。}
{$Fr<1$缓流扰动可上游，$Fr>1$急流不可上游。}

\end{multicols}
\end{document}
